% Antecedentes (consigna 2.2, ítem 6): marco normativo y estado del arte 2023-2026 (normas y fuentes de datos
% no cuentan en el estado del arte). Base: versión revisada del estudiante (25-nov-2026), con Morello (2025) en
% lugar de Larios-Hernandez (2025).
\section{Antecedentes}

\subsection{Marco normativo del financiamiento participativo financiero en el Perú}

El Decreto de Urgencia N.\textsuperscript{o}~013-2020 definió el FPF como la actividad en la que una plataforma pone en contacto a quienes buscan financiamiento, llamados receptores, con inversionistas que esperan un retorno financiero, ya sea mediante préstamos o mediante valores. Solo pueden realizarla sociedades anónimas autorizadas y supervisadas por la SMV \parencite{du0132020}. El reglamento se aprobó en mayo de 2021 \parencite{smv0452021} y se modificó en noviembre de 2025 para hacerlo más flexible. Con ese cambio subieron los límites de financiamiento por proyecto y por receptor, que quedaron en 50 UIT para los proyectos personales y en 500 UIT para los empresariales, ampliables a 1\,000 UIT si la plataforma aplica criterios de selección específicos. Además, bajo ciertas condiciones, se permitió financiar proyectos empresariales desarrollados fuera del Perú \parencite{smv0192025}. El crédito tradicional, en cambio, lo otorgan bancos y cajas que supervisa la Superintendencia de Banca, Seguros y AFP.

\subsection{Estado del arte}

\subsubsection{¿Sustituto o complemento del crédito bancario?}

Los estudios recientes no coinciden en este punto. \textcite{kim2026} modelan la elección del prestatario entre banco y plataforma. Con la entrada escalonada de LendingClub en Estados Unidos, encuentran que disminuyen los préstamos bancarios a pequeñas empresas, sobre todo en zonas de ingreso bajo. \textcite{bennaceur2026}, con más de diez mil instituciones de 57 países, hallan que la presencia fintech reduce la rentabilidad de las instituciones tradicionales. Otros trabajos van en sentido contrario. En Francia, el crédito bancario de una empresa aumenta después de un préstamo fintech, porque los activos financiados sirven de garantía \parencite{beaumont2026}, y el crédito entre pares crece más donde la banca regional atiende menos \parencite{eid2024}. Otro factor es la estructura bancaria. En China, los prestatarios de mercados bancarios concentrados que recurren a plataformas tienen mejor calidad crediticia \parencite{ma2023}, y en el modelo de \textcite{vives2025} una mayor concentración bancaria implica más volumen fintech.

\subsubsection{Precio, plazo y segmentación}

La teoría no predice que el crédito de plataforma sea siempre más barato. Según \textcite{chu2024}, el crédito fintech reduce en promedio las tasas esperadas, aunque puede empeorar el acceso de los buenos prestatarios. Para \textcite{vives2025}, la tasa fintech puede ser mayor o menor que la bancaria, según su costo de fondeo. \textcite{fareed2025} no hallan cambios en las tasas en países en desarrollo, pese a que las plataformas atienden a más prestatarios. Sobre el plazo, \textcite{serfes2025} muestran que las fintech se especializan en crédito sin garantía de corto plazo, mientras que la banca lo hace en crédito con garantía de largo plazo. La ventaja informativa de las fintech también tiene efectos reales. \textcite{hau2024} estudian las líneas de crédito automáticas que reciben los vendedores de la plataforma de Alibaba y encuentran que el crédito fintech aumenta las ventas y la satisfacción de los clientes, sobre todo en los vendedores con más asimetría de información y menos garantías.

\subsubsection{La plataforma como intermediario y su desempeño}

\textcite{braggion2025} explican por qué las plataformas de deuda adoptan rasgos de intermediario. Al hacerlo aumentan el excedente del prestamista, el beneficio de la plataforma y el crédito, aunque exponen al inversionista a riesgo de liquidez. En Italia, los bancos son más eficientes que las fintech de crédito \parencite{onorato2024}, y la eficiencia de estas últimas depende sobre todo de la rentabilidad sobre activos y de la relación entre costos e ingresos \parencite{pampurini2024}. \textcite{fasano2026} predicen los retrasos de pago en una plataforma italiana de préstamos inmobiliarios a partir de la selección adversa y la señalización. Por su parte, \textcite{miglo2025} sostiene, con un modelo de crowdfunding de recompensa, que la información asimétrica deja a este financiamiento al final del orden de preferencia de las fuentes, lo que explicaría sus volúmenes pequeños.

\subsubsection{Regulación y acceso}

Cuando las normas son claras y explícitas, aumentan el volumen del crowdfunding de deuda y la creación de plataformas \parencite{ran2025}. Los límites también cambian a quién se financia. En el crowdfunding de capital del Reino Unido, los topes para inversionistas no acreditados favorecen a empresas de regiones desfavorecidas \parencite{lazos2025}. En Estados Unidos, un fallo sobre préstamos por encima del límite de usura redujo el crédito de plataforma y aumentó las quiebras personales de hogares de bajos ingresos \parencite{danisewicz2023}. En la Unión Europea, el reglamento de 2020 excluye los préstamos a consumidores y protege al inversionista no sofisticado con técnicas del derecho del consumidor, como una prueba de conocimientos y cuatro días de reflexión, aunque esa protección es incompleta porque un resultado desfavorable no impide invertir \parencite{morello2025}. En América Latina, la inclusión financiera digital depende de la edad, la educación, el género y el ingreso \parencite{alvarezfranco2026}.

\subsubsection{Estudios sobre el Perú y vacío de conocimiento}

Sobre el Perú hay pocos trabajos y son descriptivos. \textcite{mejiatrujillo2023} caracterizan la industria fintech peruana, y \textcite{lopezaquise2026} discute si el espacio de prueba regulatorio favorece la innovación o se vuelve una barrera. Las cifras del FPF que circulan son agregados anuales de la SMV citados por la prensa \parencite{ramirez2024,ramirez2026}. La búsqueda bibliográfica para este estudio se hizo en OpenAlex y Crossref, y no se encontró ningún trabajo que caracterice la oferta de FPF peruana con datos oficiales por plataforma y mes ni que la ubique frente al crédito tradicional. Ese es el vacío que aborda este estudio.
